\documentclass[11pt]{article}
\usepackage[T1]{fontenc}
\usepackage{lmodern,amsmath,amssymb,amsthm,mathtools,mathrsfs,geometry,hyperref}
\geometry{margin=1in}
\newtheorem{theorem}{Theorem}
\newtheorem{lemma}[theorem]{Lemma}
\newtheorem{proposition}[theorem]{Proposition}
\newtheorem{corollary}[theorem]{Corollary}
\newcommand{\CC}{\mathbb C}
\renewcommand{\AA}{\mathbb A}
\newcommand{\QQl}{\overline{\mathbb Q}_\ell}
\newcommand{\End}{\operatorname{End}}
\newcommand{\Sym}{\operatorname{Sym}}
\newcommand{\im}{\operatorname{im}}
\newcommand{\coker}{\operatorname{coker}}
\newcommand{\rank}{\operatorname{rank}}
\newcommand{\tr}{\operatorname{tr}}
\newcommand{\diag}{\operatorname{diag}}
\newcommand{\Res}{\operatorname{Res}}
\title{The marked exponential family and its original arithmetic receivers\\
\large Complete receiving proof of the marked-image continuation}
\author{Split-Zero programme}
\date{19 September 2026}
\begin{document}
\maketitle

This note receives the complete nine-page \emph{Deligne marked image} calculation
\cite{received}, including its finite-field endomorphism cohomology, its tensor
constituents, and its characteristic-zero marked frame. Every formula used below
is proved with the specified full polynomial and the original coefficient frame.
The additional receiving calculations are the exact second fundamental form of
the retained polynomial line, an explicit inverse to the regularized marked
frame, and the transport of every original restriction and attained quotient
through that frame. These are finite identities. Their statements contain no
unevaluated asymptotic hypothesis.

\section{Objects, coefficient models, and imported theorems}
Retain the stipulated simple quartet with
\[
 \rho=\tfrac12+\delta+i\gamma,\qquad 0<\delta<\tfrac12,\quad\gamma>2,
 \qquad k\equiv1\pmod4,
\]
and set
\begin{equation}\tag{MII1}\label{mii:objects}
 \begin{gathered}
 c=k/2,\qquad q=(k+1)^2,\qquad
 \beta_{ab}=c+(2a-k)\delta+i(2b-k)\gamma,\\
 \chi(S)=\prod_{a,b=0}^k(S-\beta_{ab}),\qquad
 \Phi'(S)=\chi(S),\quad\Phi(0)=0,\qquad\lambda=k\rho,\\
 p_\lambda(S)=\frac{\chi(S)}{S-\lambda},\qquad
 b_\lambda=\frac1{\chi'(\lambda)},\qquad v_\lambda=b_\lambda p_\lambda.
 \end{gathered}
\end{equation}
All the roots are distinct: equality of two roots first equates their real
parts and then their imaginary parts. The numbers $\delta,\gamma$ are nonzero.
In particular $b_\lambda\ne0$. The arithmetic unit $v_h=2\xi/h$, its factor
$v_h(\rho)^k$ in the original inclusion, and the original convolution measure
remain unchanged. If that original inclusion specifies a scalar multiple
$a v_\lambda$, every occurrence of $b_\lambda$ in a frame determinant is replaced
by $a b_\lambda$; the present formulas use the cardinal polynomial in
\eqref{mii:objects} and leave the inclusion itself in its original coordinates.

The marked complex is
\begin{equation}\tag{MII2}\label{mii:complex}
 \left[\CC[u,t,S]\xrightarrow{\ u\partial_S+\chi(S)-t\ }
                 \CC[u,t,S]\,dS\right].
\end{equation}
The mark $t$, the period $u$, and the arithmetic cutoff $N$ are three separate
specified parameters. The maps connecting their uses will be given explicitly.

For the finite-field calculation choose a finitely generated coefficient ring
containing the finitely many coefficients of $\Phi$, invert $u$ and
$1,2,\ldots,q+1$, and choose a good closed fibre of characteristic $p>q+1$.
Use $\mathbb F_b$ for its finite ground field, $Q=b^n$ for an extension
cardinality, a nontrivial character $\psi:\mathbb F_p\to\QQl^*$, and $\ell\ne p$.
No reduction of the analytic function $\xi$ is used. The finite-field
polynomial and the complex polynomial arise from that coefficient model; an
identification of their cohomology fibres is not inserted.

We use the following results of Pierre Deligne with their exact scope
\cite[\S3.7, (3.7.2.3), Lemma (3.7.3); (3.4.1)(iii)]{deligne}.
For a polynomial of degree $d$ prime to $p$ whose highest homogeneous part
defines a smooth projective hypersurface, exponential compact cohomology in
$n$ variables is concentrated in degree $n$, has dimension $(d-1)^n$, and is
pure of weight $n$. In the corresponding open coefficient family these
cohomology sheaves are lisse. A lisse pointwise-pure sheaf on a normal
finite-field scheme is geometrically semisimple. For the one-variable phases
here, the highest projective zero locus is empty in $\mathbb P^0$, and the
leading coefficient is invertible. Thus all mark values are included. These
are the cited external theorems; the endomorphism computations below are
proved here. Compact base change, K\"unneth and smooth Poincar\'e duality are
used with the conventions of \cite[\S1.4]{deligne}.

\section{The entire endomorphism cohomology on the mark line}
Let
\begin{equation}\tag{MII3}\label{mii:V}
 \mathcal V=R^1\pi_!\mathcal L_\psi((\Phi(x)-tx)/u),
 \qquad\pi:\AA_x^1\times\AA_t^1\longrightarrow\AA_t^1.
\end{equation}
It is lisse of rank $q$, pure of weight one, and the full compact direct image
is $\mathcal V[-1]$.

\begin{lemma}[The exact dual and Fourier kernels]\label{mii:fourierlemma}
If $\mathcal V_-$ is defined by the negative of the complete phase, then
\begin{equation}\tag{MII4}\label{mii:dual}
 \mathcal V_-\simeq\mathcal V^\vee(-1).
\end{equation}
For $\operatorname{pr}_a:\AA_a^1\times\AA_z^1\to\AA_a^1$ and the zero
inclusion $i_0$, one has
\begin{equation}\tag{MII5}\label{mii:fourier}
 R\operatorname{pr}_{a!}\mathcal L_\psi(az)
       \simeq i_{0*}\QQl(-1)[-2].
\end{equation}
\end{lemma}
\begin{proof}
The degree $q+1$ of the phase is prime to $p$. A polar expression $g^p-g$
has highest pole order divisible by $p$, so the rank-one Artin--Schreier
character at infinity is nontrivial on wild inertia. Its wild-inertia
invariants vanish. Wild invariants are exact for $\ell$-adic coefficients
when $\ell\ne p$: one averages over the finite wild quotients, whose orders
are units. The tame-inertia spectral sequence consequently has zero input.
For the compactification $j:\AA^1\to\mathbb P^1$ this proves
$j_!\mathcal L\simeq Rj_*\mathcal L$. Poincar\'e duality pairs the two
opposite exponential degree-one groups perfectly into $\QQl(-1)$, proving
\eqref{mii:dual}, including its twist.

For \eqref{mii:fourier}, compact base change gives at $a=0$ the compact
complex of $\AA^1$, namely $\QQl(-1)[-2]$. If $a\ne0$, the cover
$y^p-y=az$ is an affine line with coordinate $y$. Its compact cohomology is
one Tate line in degree two, on which all deck translations act trivially.
The nontrivial character summand defining $\mathcal L_\psi(az)$ therefore
has zero compact cohomology. The relative complex is supported on the
closed zero point. Recollement identifies it with the direct image of its
zero-fibre complex, proving the stated relative identity.
\end{proof}

\begin{theorem}[Complete mark-line endomorphism complex]\label{mii:irreducible}
Over the algebraic closure of the specified finite ground field,
\begin{equation}\tag{MII6}\label{mii:endV}
 R\Gamma_c(\AA_t^1,\End\mathcal V)=\QQl(-1)[-2].
\end{equation}
The sheaf $\mathcal V$ is geometrically irreducible and
$H^1(\AA^1,\End\mathcal V)=0$.
\end{theorem}
\begin{proof}
Compute the compact complex of
$\mathcal L_\psi((\Phi(x)-\Phi(y)-t(x-y))/u)$ on the full ordered
space $\AA^3_{x,y,t}$. Integration first in $x,y$, the concentration in
degree one and \eqref{mii:dual} give
\[
 R\Gamma_c(\AA_t^1,\End\mathcal V)(-1)[-2].
\]
Integration first in $t$ applies \eqref{mii:fourier} to
$-(x-y)/u$. The result is supported exactly on $x=y$, with factor
$(-1)[-2]$. The remaining phase is identically zero on that diagonal, so
the same complex is
$R\Gamma_c(\AA_x^1,\QQl)(-1)[-2]$. Cancelling the identical twist and
shift proves \eqref{mii:endV}. The trace pairing identifies
$(\End\mathcal V)^\vee$ with $\End\mathcal V$. Duality on the smooth
line now gives $H^0(\End\mathcal V)=\QQl$ and $H^1(\End\mathcal V)=0$.
Deligne's geometric semisimplicity applies to the pure lisse sheaf
$\mathcal V$. In a semisimple representation $\bigoplus_i W_i^{m_i}$
over the algebraically closed coefficient field, the dimension of the
endomorphism algebra is $\sum_i m_i^2$. Its value one forces exactly one
simple summand of multiplicity one. This proves irreducibility.
\end{proof}

The corresponding exact trace identity is
\begin{equation}\tag{MII7}\label{mii:moment2}
 \sum_{t\in\mathbb F_Q}
 \left|\sum_{x\in\mathbb F_Q}
 \psi\operatorname{Tr}_{\mathbb F_Q/\mathbb F_p}
             \frac{\Phi(x)-tx}{u}\right|^2=Q^2.
\end{equation}
Indeed the $t$-sum of the expanded square is zero for $x\ne y$ and is
$Q$ for each of the $Q$ diagonal pairs. The complex identity
\eqref{mii:endV} also determines the individual cohomology groups.

A morphism from $\mathcal V$ to any lisse target has a
monodromy-invariant kernel. Thus its rank is either zero or $q$. The
same statement holds after restricting to a nonempty open subcurve: its
fundamental group surjects onto that of the original smooth curve, because
every connected finite \`etale cover remains connected on restriction to
the dense open. This is a precise restriction on a proposed lisse
extension of a fibre map. The exact characteristic-zero coefficient
and metric receivers of actual fibre maps are constructed below; no
comparison of a complex fibre with an unspecified finite-field fibre
is required by them.

\section{The complete rank-four tensor-square calculation}
Retain the one-packet centred family and primitive constant
\begin{equation}\tag{MII8}\label{mii:quintic}
 \begin{gathered}
 \phi_{\beta,\eta}(x)=x^5/5+\beta x^3/3+\eta x,\qquad
 A(\beta,\eta)=1/160+\beta/24+\eta/2,\\
 \beta=2(\gamma^2-\delta^2),\qquad
 \eta=(\delta^2+\gamma^2)^2
 \quad\hbox{at the original complex quartet}.
 \end{gathered}
\end{equation}
In characteristic $p>5$ with $u\ne0$ let
\[
 \mathcal U=R^1\pi_!\mathcal L_\psi(\phi_{\beta,\eta}(x)/u),\qquad
 \mathcal A=\mathcal L_\psi(A(\beta,\eta)/u),\qquad
 \mathcal V^{\rm orig}=\mathcal A\otimes\mathcal U.
\]
These rank-four sheaves on $\AA^2_{\beta,\eta}$ are lisse and pure of
weight one. The constant phase is retained in $\mathcal V^{\rm orig}$.

\begin{theorem}[The three tensor constituents]\label{mii:tensor}
One has
\begin{equation}\tag{MII9}\label{mii:endU}
 R\Gamma_c(\AA^2,\End\mathcal U)=\QQl(-2)[-4]
\end{equation}
and
\begin{equation}\tag{MII10}\label{mii:endtensor}
 H_c^i(\AA^2,\End(\mathcal U^{\otimes2}))=
 \begin{cases}
 \QQl&i=2,\\
 \QQl(-1)^{\oplus3}&i=3,\\
 \QQl(-2)^{\oplus3}&i=4,\\
 0&\text{otherwise}.
 \end{cases}
\end{equation}
The following direct summands are geometrically simple, pairwise
nonisomorphic, and have ranks $10,1,5$, respectively:
\begin{equation}\tag{MII11}\label{mii:decomposition}
 (\mathcal V^{\rm orig})^{\otimes2}
 =\Sym^2\mathcal V^{\rm orig}
 \oplus\mathcal A^2(-1)
 \oplus(\bigwedge^2\mathcal V^{\rm orig})_0.
\end{equation}
The last summand is the kernel of the alternating contraction to
$\mathcal A^2(-1)$.
\end{theorem}
\begin{proof}
For \eqref{mii:endU}, use ordered $x,y$. Integration in $\eta$ of the
full opposite-phase product imposes $x=y$ by \eqref{mii:fourier}.
The cubic and quintic differences then vanish. The remaining free
coordinates are $x,\beta$. The twist $(-1)$ and shift $[-2]$ from this
Fourier integration equal those from degree-one clean duality in $x,y$.
Cancelling them leaves exactly $R\Gamma_c(\AA^2,\QQl)$.

For the tensor square use ordered $x_1,x_2,y_1,y_2$. Integrating those
four fibre variables first gives
\[
 R\Gamma_c(\AA^2,\End(\mathcal U^{\otimes2}))(-2)[-4].
\]
Instead integrate $\eta,\beta$ first. Their Fourier constraints are
$x_1+x_2=y_1+y_2$ and $x_1^3+x_2^3=y_1^3+y_2^3$, with the identical
twist and shift $(-2)[-4]$. Write $x_1=x,x_2=s-x,y_1=y,y_2=s-y$.
The exact polynomial identities are
\begin{equation}\tag{MII12}\label{mii:phasefactor}
 \begin{aligned}
 x^3+(s-x)^3-y^3-(s-y)^3
     &=3s(x-y)(x+y-s),\\
 x^5+(s-x)^5-y^5-(s-y)^5
     &=5s(x-y)(x+y-s)\bigl(x^2-sx+y^2-sy+s^2\bigr).
 \end{aligned}
\end{equation}
Their signs and all the opposite-pair solutions are retained. The
second polynomial vanishes on the exact reduced zero scheme of the
first. The linear map
$(s,x,y)\mapsto(s,x-y,s-x-y)$ has determinant $-2$, so this scheme
is $Z=\{z_1z_2z_3=0\}\subset\AA^3$. Consequently
\[
 R\Gamma_c(\AA^2,\End(\mathcal U^{\otimes2}))
       =R\Gamma_c(Z,\QQl).
\]
The complement is $(\mathbb G_m)^3$. Its compact cohomology in degrees
$3,4,5,6$ is respectively
$\QQl,\QQl(-1)^3,\QQl(-2)^3,\QQl(-3)$, by three K\"unneth products
of $H_c^1(\mathbb G_m)=\QQl$ and $H_c^2(\mathbb G_m)=\QQl(-1)$.
The open--closed exact sequence in $\AA^3$ shifts the first three
groups down by one to the groups of $Z$. The top map
$H_c^6((\mathbb G_m)^3)\to H_c^6(\AA^3)$ is an isomorphism of
fundamental classes. This proves \eqref{mii:endtensor}, including its
Frobenius twists and its zero groups.

Duality on the smooth two-dimensional base gives dimensions one and
three for $H^0(\End\mathcal U)$ and
$H^0(\End(\mathcal U^{\otimes2}))$. To identify three actual
summands, the odd involution $\iota(x)=-x$ identifies the negative
phase with the positive phase. Clean degree-one duality gives
\[
 B(a,b)=\int a\smile\iota^*b,
 \qquad B:\mathcal U\otimes\mathcal U\longrightarrow\QQl(-1).
\]
It is nondegenerate. Pullback by $\iota$ preserves the top trace class
and $\iota^2=1$; exchanging the two degree-one factors changes the
sign. Therefore $B(b,a)=-B(a,b)$. In a symplectic basis
$e_1,f_1,e_2,f_2$, the invariant bivector
$e_1\wedge f_1+e_2\wedge f_2$ contracts to $2$, so it splits a
Tate line from $\bigwedge^2\mathcal U$. Its contraction kernel has
rank five and $\Sym^2\mathcal U$ has rank ten. Twisting both factors
by the original $\mathcal A$ gives \eqref{mii:decomposition}.
Deligne's semisimplicity applies to the pure lisse tensor square.
Its three nonzero direct summands contribute at least three to its
endomorphism dimension, already proved to be exactly three. Hence
each summand is simple and no two are isomorphic. Their distinct
ranks give the same last conclusion directly.
\end{proof}

For every extension field the full trace sum $S(\beta,\eta)$ satisfies
\begin{equation}\tag{MII13}\label{mii:moment4}
 \sum_{\beta,\eta\in\mathbb F_Q}|S(\beta,\eta)|^4
       =Q^2(3Q^2-3Q+1).
\end{equation}
The constant character $\mathcal A$ cancels in this fourth absolute
power. The two parameter sums give $Q^2$ times the number of points
of $Z$. At $s=0$ there are $Q^2$ pairs $(x,y)$; at each $s\ne0$
there are $2Q-1$. The exact number is
$Q^2+(Q-1)(2Q-1)=3Q^2-3Q+1$, proving the formula.

For completeness the two finite fibre maps mentioned in the received
continuation have elementary exact ranks. For
$D=\diag(\zeta,\zeta^2,\zeta^3,\zeta^4)$ with $\zeta^5=1$ primitive,
the actual cyclic average on the symmetric square is
\begin{equation}\tag{MII14}\label{mii:cyclic}
 P_5=\tfrac15\sum_{j=0}^4\Sym^2(D^j),\qquad
 \im P_5=\langle x_1x_4,x_2x_3\rangle.
\end{equation}
The characters on $x_rx_s$ are $\zeta^{r+s}$, and precisely the
two listed pairs have exponent zero modulo five. Thus the rank is
two. For the original four distinct one-packet roots $\alpha_i$,
let $e_i$ be their cardinal idempotents and let $A e_i=\alpha_i e_i$.
The nine distinct sums $\alpha_i+\alpha_j$ form the $3$ by $3$
grid with increments $2\delta,2i\gamma$. Their sum algebra maps to
the symmetric square by the exact evaluation map
\begin{equation}\tag{MII15}\label{mii:sumembedding}
 \begin{aligned}
 \iota_2:\CC[T]/\prod_{\nu}(T-\nu)&\longrightarrow\Sym^2\CC^4,\\
 f&\longmapsto f(A\otimes1+1\otimes A)(1\otimes1),\\
 \iota_2(e_\nu)&=
 \sum_{2\alpha_i=\nu}e_i\otimes e_i
 +2\sum_{i<j,\ \alpha_i+\alpha_j=\nu}e_i\odot e_j,
 \end{aligned}
\end{equation}
where $e_i\odot e_j=(e_i\otimes e_j+e_j\otimes e_i)/2$.
The nine displayed images are nonzero and have disjoint supports in
the symmetric tensor basis. This proves injectivity and rank nine,
including the multiplicity two of an off-diagonal symmetric pair.
Neither this image nor \eqref{mii:cyclic} can be the image of a
proper nonzero lisse projector on the simple rank-ten summand of
\eqref{mii:decomposition}. These exact finite maps remain available
as fibre maps; the proved assertion concerns their extension as
projectors of this particular finite-field family.

\section{The marked coefficient module and the complete derivative frame}
Return to the complex polynomial \eqref{mii:objects}, put
$R=\CC[u,u^{-1},t]$, and write
\[
 \mathscr H=R[S]/(u\partial_S+\chi-t)(R[S]).
\]
This is a quotient of $R$-modules. Polynomial division in the
differential expression gives a unique representative of degree less
than $q$: subtract the leading coefficient times
$(u\partial_S+\chi-t)S^{d-q}$ from a polynomial of degree $d\ge q$,
and repeat. Termination follows from the strictly decreasing degree.
Uniqueness follows because every nonzero polynomial $f$ has
$\deg((u\partial_S+\chi-t)f)=q+\deg f\ge q$.
Thus $1,S,\ldots,S^{q-1}$ is a free $R$-basis.

The cochain commutator and its induced connection are
\begin{equation}\tag{MII16}\label{mii:connection}
 [\partial_t-S/u,\ u\partial_S+\chi-t]=0,
 \qquad\nabla_t=\partial_t-A(t)/u.
\end{equation}
Here $A(t)$ is the companion matrix of $\chi-t$ on these
representatives: multiplying a polynomial of degree at most $q-1$
by $S$ requires only the relation
$(u\partial_S+\chi-t)1=\chi-t$. This proves the displayed matrix
connection. It does not define a ring multiplication on the
differential quotient. At $t=0$ the representative map identifies
its underlying vector space with $E=\CC[S]/\chi$ exactly.

Let $R_j$ denote the representative of $[S^j p_\lambda]$. The
polynomial relation $(S-\lambda)p_\lambda=\chi$ and reduction of
$(\chi-t)S^{j-1}$ give
\begin{equation}\tag{MII17}\label{mii:recurrence}
 R_j-\lambda R_{j-1}
       =t S^{j-1}-u(j-1)S^{j-2},\qquad 1\le j\le q.
\end{equation}
The last term is zero at $j=1$. Its right side already has degree
less than $q$, so this is equality of the literal representatives.
Applying the cochain connection to the unreduced fixed polynomial
before reduction gives, for every $j\ge0$,
\begin{equation}\tag{MII18}\label{mii:jets}
 \nabla_t^j[v_\lambda]=b_\lambda(-u)^{-j}R_j.
\end{equation}
The commutator in \eqref{mii:connection} proves that differentiating
the reduced representative gives exactly the same class.

\begin{theorem}[Determinants and the entire supported cokernel]
In the increasing original $S$-coefficient basis, let
\[
 \mathcal O_\lambda(t)=[v_\lambda,\nabla_t v_\lambda,
                  \ldots,\nabla_t^{q-1}v_\lambda],
 \qquad\epsilon_q=(-1)^{(q-1)(q+2)/2}.
\]
Then
\begin{equation}\tag{MII19}\label{mii:Odet}
 \det\mathcal O_\lambda(t)
   =\epsilon_q b_\lambda^q u^{-q(q-1)/2}t^{q-1},
 \qquad\coker\mathcal O_\lambda\simeq R/(t^{q-1}).
\end{equation}
The fibre matrix has rank $q-1$ at zero. The regularized original
frame
$\mathcal J_\lambda=[v_\lambda,\nabla_t^2v_\lambda,
\ldots,\nabla_t^qv_\lambda]_{t=0}$ is invertible with
\begin{equation}\tag{MII20}\label{mii:Jdet}
 \det\mathcal J_\lambda
    =\epsilon_q(q-1)!b_\lambda^q u^{-q(q-1)/2}.
\end{equation}
\end{theorem}
\begin{proof}
First remove the factors $b_\lambda(-u)^{-j}$ from the columns
of $\mathcal O_\lambda$. Perform the unipotent column operation
$R_j\mapsto R_j-\lambda R_{j-1}$ simultaneously, or from right
to left. Equation \eqref{mii:recurrence} gives the columns
$p_\lambda$ and $tS^{j-1}-u(j-1)S^{j-2}$, $1\le j\le q-1$.
The first is monic of degree $q-1$; the other columns have degrees
at most $q-2$. Moving the first column to the last position gives
the sign $(-1)^{q-1}$. The remaining upper-bidiagonal block has
diagonal entries $t$. Restoring the column factors gives sign
$(-1)^{q-1+q(q-1)/2}=\epsilon_q$ and proves the determinant.

The first column eliminates the top coefficient generator. On the
remaining generators $e_0,\ldots,e_{q-2}$ the presentation is exactly
\[
 te_0=0,\qquad te_j=u j e_{j-1}\quad(1\le j\le q-2).
\]
The explicit presentation isomorphism is
\begin{equation}\tag{MII21}\label{mii:cokeriso}
 e_j\longmapsto
 \frac{j!}{(q-2)!}u^{-(q-2-j)}t^{q-2-j}\pmod{t^{q-1}},
 \qquad\text{inverse: }1\longmapsto e_{q-2}.
\end{equation}
Every relation maps to zero. Conversely the displayed relations,
using the invertible constants $uj$, express every generator in
terms of $e_{q-2}$ by precisely this formula; its sole remaining
relation is $t^{q-1}e_{q-2}=0$. This proves the isomorphism and
shows that the special fibre of the cokernel has dimension one.

For the regularized frame remove the column scalars with orders
$0,2,\ldots,q$. At zero $R_1=\lambda p_\lambda$; replace
$R_2$ by $R_2-\lambda^2p_\lambda$, and, from right to left, each
$R_j$ for $j\ge3$ by $R_j-\lambda R_{j-1}$. The columns become
$p_\lambda,-u,-2uS,\ldots,-(q-1)uS^{q-2}$. Their determinant
is $u^{q-1}(q-1)!$. The restored exponent of $-u$ is
$-[q(q+1)/2-1]$, giving exactly \eqref{mii:Jdet}.
\end{proof}

The proof of this finite algebraic theorem also works for a
terminal primary class when
$\chi=(S-\lambda)^e\chi_\lambda$ and $\chi_\lambda(\lambda)\ne0$:
use $p_\lambda=\chi/(S-\lambda)$ and
$b_\lambda=1/\chi_\lambda(\lambda)$. This statement follows from
the same polynomial recurrence, which used no simple-root
division. It does not change the simple-quartet metric in this note.

\section{The explicit inverse and the retained-line leakage}
Put $v_j=(\nabla_t^j v_\lambda)|_{t=0}$. Formula
\eqref{mii:recurrence} gives the following inverse reconstruction
in the original coefficient space:
\begin{equation}\tag{MII22}\label{mii:inverse}
 \begin{aligned}
 1&=-\frac{u}{b_\lambda}v_2
            +\frac{\lambda^2}{b_\lambda u}v_0,\\
 S^r&=\frac{(-u)^{r+1}v_{r+2}
                    -\lambda(-u)^r v_{r+1}}{b_\lambda(r+1)}
                &&(1\le r\le q-2),\\
 S^{q-1}&=b_\lambda^{-1}v_0-
                         \sum_{r=0}^{q-2}p_rS^r,
 \qquad p_\lambda=S^{q-1}+\sum_{r=0}^{q-2}p_rS^r.
 \end{aligned}
\end{equation}
To verify the first identity, set $j=2,t=0$ in the recurrence to
obtain $R_2=\lambda^2p_\lambda-u$. For the middle identities use
$R_{r+2}-\lambda R_{r+1}=-u(r+1)S^r$ and then
$R_j=b_\lambda^{-1}(-u)^jv_j$. The last identity is the monic
definition of $p_\lambda$. These identities express each original
basis vector in the $q$ specified jet columns and are an explicit
inverse matrix for $\mathcal J_\lambda$; no basis existence argument
replaces the map.

The exact bundle map measuring the retained polynomial line is
also explicit. Let $L=R v_\lambda\subset\mathscr H$, and let
$\pi_L:\mathscr H\to\mathscr H/L$ be its quotient map. Since
$v_\lambda$ is monic up to the unit $b_\lambda$, $L$ is a direct
summand as an $R$-module. The second fundamental form of this
specified line is the $R$-linear map
\begin{equation}\tag{MII23}\label{mii:secondfund}
 \begin{aligned}
 \mathfrak b_L:L&\longrightarrow(\mathscr H/L)\,dt,\\
 f v_\lambda&\longmapsto
       \pi_L(\nabla_t(fv_\lambda))\,dt
       =-\frac{b_\lambda t}{u}f\,\pi_L(1)\,dt.
 \end{aligned}
\end{equation}
The Leibniz term $f'v_\lambda$ maps to zero. The formula itself
is \eqref{mii:recurrence} at $j=1$. The class $\pi_L(1)$ is
nonzero because $q\ge2$ and a degree-zero polynomial cannot be
an $R$-multiple of the nonconstant $p_\lambda$. In particular,
the second fundamental form vanishes to exactly first order
at the original mark, and
\begin{equation}\tag{MII24}\label{mii:secondjet}
 \pi_L(v_1)=0,\qquad
 \pi_L(v_2)=-\frac{b_\lambda}{u}\pi_L(1)\ne0.
\end{equation}
This computes the first actual marked leakage from the original
arithmetic eigenline $A(0)v_\lambda=\lambda v_\lambda$.
Any connection-stable subbundle containing the specified
polynomial section contains every $v_j$ and hence, by
\eqref{mii:inverse}, the entire fibre at zero.

There is also an exact formal bridge to a horizontal extension
of its initial vector. Write
$A(t)=A_0+tB$, $B=e_0e_{q-1}^{\mathsf T}$ in the original
increasing coefficient frame. The formal matrix
$H(t)=\sum_{n\ge0}H_nt^n$ is uniquely defined by
\begin{equation}\tag{MII25}\label{mii:horizontal}
 H_0=I,\quad H_{-1}=0,\qquad
 (n+1)H_{n+1}=u^{-1}(A_0H_n+B H_{n-1}).
\end{equation}
It has invertible constant coefficient and satisfies
$\nabla_t H=0$. Thus $a(t)\mapsto H(t)a(t)$ is a
connection-preserving isomorphism from the trivial formal
connection to the completed marked connection. Its line through
$v_\lambda$ is the horizontal extension of that fibre vector.
The polynomial line's actual leakage remains
\eqref{mii:secondfund}; these two sections are related by the
explicit matrix \eqref{mii:horizontal}, without identifying them.

\section{Period determinant, Gamma constants, and scalar equation}
The full root polynomial satisfies $\chi(c+w)=\chi(c-w)$.
Its degree $q$ is even. Write
\[
 \Phi(c+w)=C_\Phi+\phi(w),\quad C_\Phi=\Phi(c),\quad
 \phi(0)=0,\quad\phi\text{ odd}.
\]
Choose the original branch of $u^{1/(q+1)}$ and retain the ordered
decay rays with angles $(\arg u+\pi+2\pi j)/(q+1)$, $0\le j\le q$.
The cycles are the original differences $\Gamma_j=-\ell_0+\ell_j$,
$1\le j\le q$. Moving their common finite starting point to
$w=0$ changes no integral: the same finite connecting path cancels
between the two rays, and the joining arcs at infinity decay in
the same sectors. Define the period matrix by
\[
 \Pi_{ja}(u,t)=\int_{\Gamma_j}w^{a-1}
    \exp\bigl((C_\Phi-ct+\phi(w)-tw)/u\bigr)\,dw,
 \qquad1\le j,a\le q.
\]
Let
\begin{equation}\tag{MII26}\label{mii:Fg}
 \zeta=e^{2\pi i/(q+1)},\quad F_{ja}=\zeta^{ja}-1,\qquad
 g_a=(q+1)^{a/(q+1)-1}e^{\pi ia/(q+1)}
             \Gamma\!\left(\frac a{q+1}\right).
\end{equation}

\begin{theorem}[Complete original period determinant]\label{mii:periodtheorem}
With the just specified branches, orientations and column order,
\begin{equation}\tag{MII27}\label{mii:perioddet}
 \det\Pi(u,t)=
 e^{q(C_\Phi-ct)/u}(\det F)
                 \left(\prod_{a=1}^qg_a\right)u^{q/2}.
\end{equation}
\end{theorem}
\begin{proof}
Separate the constant $C_\Phi-ct$ in the exponent. Vary a lower
odd coefficient of the centred phase, with inserted odd power
$w^p$, $p\le q-1$. For the column $w^j$, $0\le j\le q-1$,
perform ordinary division
$w^{p+j}=(\chi(c+w)-t)Q(w)+r(w)$. If the quotient is nonzero
its degree is at most $p+j-q\le j-1$. In the differential
quotient, the inserted column is represented by
$r-uQ'$. This last term has degree at most $j-2$ and needs no
further reduction. The ordinary remainder map is multiplication
by $w^p$ in the quotient by the even polynomial $\chi(c+w)-t$;
it reverses parity and so has zero trace. The correction
$-uQ'$ is strictly lower triangular and also has zero trace.
Therefore the coefficient derivative of the period matrix has
the form $u^{-1}\Pi M_p$ with $\tr M_p=0$. Multilinearity, or
the adjugate formula, gives zero derivative of its determinant
without presuming its invertibility. This proves independence
of every lower odd coefficient, including the marked coefficient
$-t$ in the centred phase.

All these differentiations are dominated along the fixed rays
on compact coefficient sets by
$\exp[-R^{q+1}/((q+1)|u|)+O(R^{q-1})]$ times a polynomial.
At the monomial phase, use
$w=u^{1/(q+1)}e^{\pi i/(q+1)}\zeta^j r$ with $r>0$ on the
$j$th ray, followed by $s=r^{q+1}/(q+1)$. The literal integral
is $\zeta^{ja}u^{a/(q+1)}g_a$. Taking the specified differences
gives $F\diag(u^{a/(q+1)}g_a)$. Finally
$\sum_{a=1}^q a/(q+1)=q/2$, and restoring the separated
constant proves \eqref{mii:perioddet}.
\end{proof}

Let $\mathbf y_\lambda(t)$ be the vector of integrals of the
original amplitude $v_\lambda(S)$. Translation of coefficient
vectors from $S$ to $w$ is the exact matrix
\[
 T_{rj}=\binom jr c^{j-r}\quad(0\le r\le j\le q-1),\qquad
 T_{rj}=0\quad(r>j),\qquad\det T=1.
\]
Differentiating integrals gives the cochain connection in
\eqref{mii:connection}. Hence the derivative matrix is literally
$\Pi T\mathcal O_\lambda$, and
\begin{equation}\tag{MII28}\label{mii:wronskian}
 \det[\mathbf y_\lambda,\mathbf y'_\lambda,\ldots,
                       \mathbf y^{(q-1)}_\lambda]
 =\epsilon_q b_\lambda^q u^{-q(q-1)/2}t^{q-1}\det\Pi(u,t).
\end{equation}
The regularized determinant at the original mark is
\begin{equation}\tag{MII29}\label{mii:regperiod}
 \left|\det[\mathbf y_\lambda,\mathbf y''_\lambda,\ldots,
                    \mathbf y_\lambda^{(q)}]_{t=0}\right|^2
 =(2\pi)^q((q-1)!)^2|b_\lambda|^{2q}
     |u|^{-q(q-2)}e^{2q\Re(C_\Phi/u)}.
\end{equation}
To check every constant, summing characters gives
$F^*F=(q+1)(I+\mathbf1\mathbf1^*)$, whence
$|\det F|^2=(q+1)^{q+1}$. Gauss's multiplication formula
\cite[(5.5.6)--(5.5.7)]{dlmf} gives
$\prod_{a=1}^q\Gamma(a/(q+1))=(2\pi)^{q/2}/\sqrt{q+1}$.
Consequently $|\det F\prod g_a|^2=(2\pi)^q$.
Combining this with \eqref{mii:Jdet} and
\eqref{mii:perioddet} proves \eqref{mii:regperiod}.

For any one original cycle let $Y$ be its period with amplitude
$v_\lambda$, let $F_0$ be its period with amplitude one, and
put $D=-u\partial_t$. Integration by parts in $S$ gives
$(D-\lambda)Y=b_\lambda tF_0$ and
$Y=b_\lambda p_\lambda(D)F_0$. For $t\ne0$,
$tDt^{-1}=D+u/t$, so
\begin{equation}\tag{MII30}\label{mii:scalar}
 [p_\lambda(D+u/t)(D-\lambda)-t]Y=0.
\end{equation}
The operator factors have this order. The leading two coefficients
are $D^q+[\chi_{q-1}+(q-1)u/t]D^{q-1}$, where
$\chi_{q-1}=-qc$. Abel's determinant equation therefore gives
$W'/W=(q-1)/t-qc/u$, as in \eqref{mii:wronskian}.
Every component period is holomorphic at zero. Its apparent
singularity is represented by the explicit cyclic cokernel
\eqref{mii:cokeriso} and the regularized frame, without discarding
its multiplicity $q-1$.

\section{Exact receiving maps for the original arithmetic metrics}
Fix $u\ne0$ in the original branch. Let $G_N$ be the unchanged
positive Hermitian metric on the original $q$-dimensional space
$E=\CC[S]/\chi$, using its actual masses, cutoffs, root factors,
and full arithmetic unit. Set $J=\mathcal J_\lambda$ with its
explicit inverse \eqref{mii:inverse}. The diagram of vector
spaces and metrics is the literal pullback
\begin{equation}\tag{MII31}\label{mii:gram}
 J:(\CC^q,J^*G_NJ)\longrightarrow(E,G_N),\qquad
 \det(J^*G_NJ)=((q-1)!)^2|b_\lambda|^{2q}
                       |u|^{-q(q-1)}\det G_N.
\end{equation}
The map $J$ sends standard vectors to the original specified
jet columns; it introduces no source normalization. The determinant
identity follows by multiplying determinants and using
\eqref{mii:Jdet}. For the actual positive extremal root the
normalizing constant itself is completely specified:
\begin{equation}\tag{MII32}\label{mii:rootproduct}
 |b_\lambda|^{-1}
  =\prod_{\substack{0\le i,j\le k\\(i,j)\ne(0,0)}}
                   2\sqrt{\delta^2 i^2+\gamma^2j^2}.
\end{equation}
This is the product of all differences
$\lambda-\beta_{ab}=2(k-a)\delta+2i(k-b)\gamma$.

The map also transports every specified arithmetic subquotient,
not only its total determinant. Let $I_X:\CC^r\hookrightarrow E$
be any actual injective column map and $Q_X:E\twoheadrightarrow
\CC^s$ any actual surjection. Define their jet-coordinate maps
$\widetilde I_X=J^{-1}I_X$ and $\widetilde Q_X=Q_XJ$.
Then
\begin{equation}\tag{MII33}\label{mii:restrictions}
 \begin{aligned}
 \widetilde I_X^*(J^*G_NJ)\widetilde I_X&=I_X^*G_NI_X,\\
 \widetilde Q_X(J^*G_NJ)^{-1}\widetilde Q_X^*
                         &=Q_XG_N^{-1}Q_X^*,\\
 \left[\widetilde Q_X(J^*G_NJ)^{-1}\widetilde Q_X^*\right]^{-1}
                         &=(Q_XG_N^{-1}Q_X^*)^{-1}.
 \end{aligned}
\end{equation}
The first equality cancels $J,J^{-1}$ directly. For the second
use $(J^*G_NJ)^{-1}=J^{-1}G_N^{-1}J^{-*}$; the third follows
because surjectivity and positivity make both matrices invertible.
Equivalently the change $x=Jz$ bijects the exact affine fibres
$\{x:Q_Xx=y\}$ and $\{z:\widetilde Q_Xz=y\}$ and preserves their
quadratic values. Their attained minima agree. Restricting first
and taking a further quotient applies these same equalities to
$I_X^*G_NI_X$ and the unchanged next quotient map. Thus every
finite composition of the programme's actual restrictions and
attained quotients has this exact receiver. No rank-proportional
allocation has been substituted for any of these matrices.

Retain the original four-endpoint signs by defining
\[
 \mathcal R_{\rm ar}f=f(q-1)+f(q)-f(2q-1)-f(2q).
\]
The multiplier in \eqref{mii:gram} is independent of $N$;
$1+1-1-1=0$. Therefore
\begin{equation}\tag{MII34}\label{mii:return}
 \mathcal R_{\rm ar}\log\det(J^*G_NJ)
    =\mathcal R_{\rm ar}\log\det G_N
    =\mathcal B_k^{\rm ar}.
\end{equation}
For a restriction or an attained quotient transported by
\eqref{mii:restrictions}, even its full matrix is identical in its
original domain or codomain, so its signed return agrees term by
term. This supplies the explicit morphisms relating the geometric
marked frame to the arithmetic receivers. It does not evaluate
the still-present matrices $I_X^*G_NI_X$ or $Q_XG_N^{-1}Q_X^*$.
In particular the seven original allocations and the three-column
current responses retain those exact matrices. Their calculation
continues from these equalities with the original inverse and
complete action retained.

\section{Acceptance scope and provenance}
The admitted conclusions are the full formulas MII6--MII13 for
the specified good finite-field families, MII16--MII30 for the
original complex marked family, and the exact original-metric
receivers MII31--MII34. MII22--MII25 and MII33 explicitly extend
the received calculation. They retain the original full root
polynomial, mark, period branch, and all coefficient factors.
The degree-two ranks in MII14--MII15 were proved here independently.

The received text also quotes a large-$k$ five-orbit observation
rank $q-(8k-16)$. Its precise admission domain and its earlier
proof are not an input to any theorem in this receiving note;
the formula is not newly certified here. The general obstruction
to an intermediate-rank lisse morphism is proved above without it.
No assertion here identifies an arithmetic Frobenius, makes a
cross-characteristic fibre isomorphism, or evaluates an absolute
native angular allocation. In particular none asserts that the
broader tau-base or $\mathbb F_1$ programme has failed.

The entire received PDF and transcript were read. The local full
D032 source was consulted specifically at (3.4.1)(iii) and
(3.7.2.3) with its following lissity argument. The primary Numdam
record and the DLMF multiplication formula were checked on
19 September 2026. The source hashes, exact scope, and independent
algebraic checks are recorded in the accompanying machine-readable
acceptance file. Finite rational fixtures are used only to detect
algebraic errors; they do not represent actual zeros of $\xi$.

\begin{thebibliography}{9}
\bibitem{deligne}
Pierre Deligne, \emph{La conjecture de Weil. II}, Publications
Math\'ematiques de l'IH\'ES \textbf{52} (1980), 137--252.
\href{https://doi.org/10.1007/BF02684780}{doi:10.1007/BF02684780}.
Primary record:
\url{https://www.numdam.org/item/PMIHES_1980__52__137_0/}.
Used: (1.4), (3.4.1)(iii), (3.7.2.3), and the lissity assertion
in (3.7.3). Local supplied full English mathematical source:
\texttt{D032\_FULL\_EN\_source\_checked\_corrected.tex}.

\bibitem{dlmf}
NIST Digital Library of Mathematical Functions,
\emph{Gamma Function}, \S5.5(iii), Gauss's multiplication formula,
equations (5.5.6) and (5.5.7).
\url{https://dlmf.nist.gov/5.5.E6};
\url{https://dlmf.nist.gov/5.5.E7}.

\bibitem{received}
Split-Zero programme, \emph{The original marked exponential family:
Deligne constituents and the retained-class jet divisor},
19 September 2026, nine-page received PDF
\texttt{DELIGNE\_MARKED\_IMAGE.pdf}, DMI1--DMI27.
The accompanying user-supplied transcript contains the same
calculation and its input/output record. Its source SHA256 and
the exact local paths are preserved in \texttt{ACCEPTANCE.json}.
This reference credits the received calculation separately from
the human results explicitly cited above.
\end{thebibliography}
\end{document}
